\scanpage{28}
5. La donnée du \deleted{\unsure{$\Omega$}} rev. quadratique $\Omega$ sur $S$,
grâce à la considération du plongmt can. de
$\Omega$ dans son enveloppe affine $E$, \unsure{revient|vient} à la donnée

a) d'un fibré affine de rang 1 $E$ sur $S$

b) d'un diviseur relatif \deleted{de} $\Omega$ de $E/S$,
i.e. d'une section de $\underline{\mathrm{Div}}^{+2}_{E/S}\simeq\Sigma^2(E)=E\times_S E/\unsure{\mathfrak S_2|\mathscr S_2}$.

On sait que \deleted{\unsure{les}} les $\underline{\mathrm{Div}}^{\unsure{+n|n}}_{E/S}$ \inserted{$\simeq(E/S)^n/\mathfrak S_n$}
sont des fibrés affines sur $S$, de rangs respectifs $n$, qu'on peut
obtenir ainsi. On regarde d'abord le cas où $E$
est trivialisé par une section $s$
\[
E\overset{\alpha_s}{\xleftarrow{\sim}}\check{\mathbb V}L \tag*{(27)}
\]
et on a alors un iso
\[
\Sigma^n(\alpha_s):\ \Sigma^n\check{\mathbb V}L)=\check{\mathbb V}\left(\sum_{1\le i\le n}L^{\otimes i}\right)\xrightarrow{\sim}\Sigma^n E
\]
où l'hom
\[
\Sigma^n\check{\mathbb V}L\simeq\check{\mathbb V}L^n/\mathfrak S_n\longrightarrow
 \check{\mathbb V}\left(\sum_{1\le i\le n}L^{\otimes i}\right) \tag*{(28)}
\]
est défini par passage au quotient à partir de
\[
\sigma_{\unsure{n|\bullet}}:\check{\mathbb V}(L^n)\longrightarrow
 \check{\mathbb V}\Bigl(\sum_{1\le i\le n}L^{\otimes i} \tag*{(29)}
\]
\unsure{donnés|donné} par
\[
\sigma_{\unsure{\bullet|n}}=(\sigma_i)_{1\le i\le n}\qquad
\sigma_i(\deleted{\unsure{\zeta|x}}u_1,\text{---},\deleted{\unsure{\zeta|x}}u_n)=\sum u_1 u_2\text{---}u_i \tag*{(30)}
\]
($i$\textsuperscript{ème} fonction symétrique élémentaire) On aura

si
\[
s'=s+\tau\qquad(\tau\in\Gamma L)
\]
i.e. \qquad $\deleted{\unsure{s}}\alpha_{s'}=\alpha_s+\operatorname{trans}(\tau)$ \qquad d'où

% archive sheet number: 27

